\scanpage{402}\mspage{141}
ce qui donnera aussi
\[
\overline{c^{\mathrm{opp}}}=\bar c^{-1}
\]
(en admettant que les chemins triviaux $a\overset{\mathrm{id}}{\unsure{\dashv|\to}}a$
donnent l'identité $\overline{X'}_a\overset{\substack{\mathrm{id}\\[-2pt]\sim}}{\relbar}X'_a$ !]. \deleted{\unsure{En d'autre}}

Bien sûr, $\bar c$ ne doit dépendre que de la
\uline{classe d'homotopie} de $\bar c$ (dans $\underline{\mathrm{Ch}\deleted{\unsure{\mathrm{loc}}}}_{\tau}(X;a,b)$ \deleted{\unsure{ou}}
\deleted{\unsure{ou mieux,}} \inserted{ou même} \deleted{dans} $C=\underline{\unsure{\mathrm{Ch}|\mathrm{ChP}}}(X,a,b)$, catégorie de
tous les chemins dans $X$, de $a$ vers $b$, de
type non précisé). En d'autres termes, on
doit avoir un foncteur
\[
\begin{aligned}
\Pi(X)&\longrightarrow\mathrm{Hot}\\
a&\unsure{\mathrel{\vdash\!\dashv}|\longmapsto}\mathrm{hot}(\overline{X'}_a)
\end{aligned}
\]
du groupoïde fondamental de $X$ dans $\mathrm{Hot}$.
\deleted{\unsure{Et on aura}} À vrai dire, ceci n'a rien à
voir avec la situation particulière étudiée
ici, -- ça \deleted{\unsure{dev}}doit marcher chaque fois qu'on
a une $\mathrm{HOT}$-fibration, ici $\overline{X'}\longrightarrow X$,
% slanted
\margin{Ici aussi, il faudrait décrire l'effet du transport $\bar c$ sur le $\unsure{\bar\pi_0|\pi_0}$ des fibres.}
mais aussi $\underline{\mathrm{Ch}}_{\infty}(X)\longrightarrow X\times X$\textsuperscript{*}. Et
pour une $\mathbb D$-fibration \inserted{$X'\to X$} ($\mathbb D$ un dérivateur
quelconque), \inserted{et $a\in\mathbb D(e)$,} on doit avoir de $=$ deux foncteurs
\[
\Pi(X)\rightrightarrows\mathbb D(e)
\]
correspondants à
\[
\left\{\begin{aligned}
x&\longmapsto H^{\bullet}_{\mathbb D}(X'_{\unsure{x|*}},a)\\
\unsure{x|\lambda}&\unsure{\mapstochar\relbar|\mapsto}H^{\mathbb D}_{\bullet}(X'_x,a)
\end{aligned}\right.
\]
% archive sheet number: 401
